% Sentences that begin on pp.~47--48.
% The proof of Proposition~6.8.1 begins on p.~46. The sentence that
% ends there (``fully faithful \ldots'') is not repeated here.

By hypothesis, one has an isomorphism
\(\theta\colon\mathrm{pr}_2^*(\mathcal{F})\isomto\mu^*(\mathcal{F})\);
taking the inverse image of~\(\theta\) by the morphism
\(\tau\colon G\to G\times_S G\) with components \((\mathrm{id}_G,\varepsilon\circ\pi)\),
one obtains an isomorphism \(\mathcal{F}\isomto\pi^*\varepsilon^*(\mathcal{F})\).

\begin{remarks}{6.8.2}\label{I.6.8.2}
(a)~Let us consider the action of \(G\times_S G\) on~\(G\) defined by
\((g_1,g_2)\cdot g=g_1\,g\,g_2^{-1}\);
then the stabilizer of the unit section \(\varepsilon\colon S\to G\) is the
diagonal subgroup~\(H\) of \(G\times_S G\). Consequently, if~\(\mathcal{F}\)
is a quasi-coherent \((G\times_S G)\)-equivariant \(\OO_G\)-module, then,
by~6.5.2, \(\mathcal{E}=\varepsilon^*(\mathcal{F})\) is endowed with the
structure of an \(H\)-\(\OS\)-module.

(b)~One can show that \(\mathcal{F}\mapsto\varepsilon^*(\mathcal{F})\) is an
equivalence of categories, between the category of quasi-coherent
\((G\times_S G)\)-equivariant \(\OO_G\)-modules and that of quasi-coherent
\(H\)-\(\OS\)-modules. This is a particular case of more general ``descent''
results (\cf Exp.~IV, \S 2 and SGA~1, VIII); see for example \cite{Th87},
1.2--1.3.
\end{remarks}

\begin{remark}{6.8.3}\label{I.6.8.3}
One keeps the notations of~6.8.1. For every quasi-coherent
\(\OS\)-module~\(\mathcal{E}\), let us denote by \(\pi_*^G\,\pi^*(\mathcal{E})\)
the sub-\(\OS\)-module of \(\pi_*\pi^*(\mathcal{E})\) whose sections over
every open~\(V\) of~\(S\) are the
\(\gamma\in\Gamma(\pi^{-1}(V),\pi^*(\mathcal{E}))\) such that
\(g\cdot\gamma_{S'}=\gamma_{S'}\) for every \(S'\to V\) and every
\(g\in G(S')\). Then the natural morphism
\(\mathcal{E}\to\pi_*^G\,\pi^*(\mathcal{E})\) is an isomorphism: this is
immediate if \(\pi_*\pi^*(\mathcal{E})=\mathcal{E}\otimes_{\OS}\pi_*(\OO_G)\)
(for example if \(G\to S\) is affine, or if \(G\to S\) is quasi-compact and
quasi-separated and~\(\mathcal{E}\) is flat), and this is verified without
difficulty in the general case.
\end{remark}

\begin{remark}{6.8.4}\label{I.6.8.4}
Let \(S\) be a scheme, \(H\) an \(S\)-group scheme operating on an
\(S\)-scheme~\(X\), and \(\mathcal{F}\) an \(\OX\)-module. One assumes
that~\(H\) is \emph{flat} over~\(S\), and one denotes by
\(\mathbf{W}_{\mathrm{P}}(\mathcal{F})\) the restriction of the functor
\(\mathbf{W}(\mathcal{F})\) to the full subcategory consisting of the
\emph{flat} \(S\)-schemes.
Since, by~6.5.1, endowing~\(\mathcal{F}\) with an \(H\)-equivariant module
structure is equivalent to giving an isomorphism
\(\theta\colon\mathrm{pr}_X^*(\mathcal{F})\to\lambda^*(\mathcal{F})\) of
sheaves on \(G\times_S X\), satisfying the ``cocycle condition''
\((\star\star)\), one sees that, in order to give an \(H\)-equivariant
module structure on~\(\mathcal{F}\), it suffices to give such a structure
on \(\mathbf{W}_{\mathrm{P}}(\mathcal{F})\).
% typo: cocyle -> cocycle
% typo?: kept as printed (sheaves on \(G\times_S X\), while the group is \(H\))
\end{remark}

\begin{enonce}{Recall}{6.8.5}\label{I.6.8.5}
Let \(X\) be an \(S\)-scheme and \(Y\) a sub-\(S\)-scheme of~\(X\). One denotes
by \(\mathcal{N}_{Y/X}\) the conormal sheaf of the immersion
\(i\colon Y\hookrightarrow X\) (\cf EGA~IV\(_4\), 16.1.2). If \(S'\to S\) is a
\emph{flat} morphism and if one denotes by \(i'\colon Y'\hookrightarrow X'\)
the immersion deduced from~\(i\) by base change, then, by \loccit,
16.2.2~(iii), one has
\(\mathcal{N}_{Y/X}\otimes_{\OO_Y}\OO_{Y'}=\mathcal{N}_{Y'/X'}\).
\end{enonce}

\begin{proposition}{6.8.6}\label{I.6.8.6}
Let \(S\) be a scheme, \(X\) an \(S\)-group scheme, and \(Y\) a group
sub-\(S\)-scheme of~\(X\). One assumes that~\(Y\) is \emph{flat} over~\(S\).
Then the conormal sheaf \(\mathcal{N}_{Y/X}\) is a
\((Y\times_S Y)\)-equivariant \(\OO_Y\)-module.
\end{proposition}

Indeed, \(H=Y\times_S Y\) is flat over~\(S\), so, by remark~6.8.4, it
suffices to endow \(\mathbf{W}_{\mathrm{P}}(\mathcal{N}_{Y/X})\) with the
structure of an \(H\)-equivariant module. Let \(S'\) be a flat \(S\)-scheme,
let \(Y'\hookrightarrow X'\) be the immersion obtained by base change, and
let \(\mathcal{N}'=\mathcal{N}_{X/Y}\otimes_{\OO_Y}\OO_{Y'}\).
% typo?: kept as printed (\(\mathcal{N}_{X/Y}\); 6.8.5 has \(\mathcal{N}_{Y/X}\))
By~6.8.5, one has \(\mathcal{N}'=\mathcal{N}_{Y'/X'}\).

Every \(h\in H(S')\) induces an automorphism of~\(Y'\), and one therefore
obtains, for every \(y\in Y'(S')\), isomorphisms
\[
\Gamma(S',y^*(\mathcal{N}'))\isomto\Gamma(S',y^*h^*(\mathcal{N}'))
\]
which endow \(\mathbf{W}_{\mathrm{P}}(\mathcal{N})\) with the structure of an
\(H\)-equivariant module (\cf 6.1).

% Printed as a heading on the same page, not as a new chapter.
\makeatletter
\renewcommand{\@bibtitlestyle}{%
  \par\addvspace{1.1\baselineskip}%
  \noindent{\large\bfseries Bibliography}%
  \nde{: additional references cited in this Exposé.}%
  \par\addvspace{0.45\baselineskip}%
}
\makeatother
\begin{thebibliography}{DG70}

\bibitem[DG70]{DG70}
M.~Demazure, P.~Gabriel, \emph{Groupes Algébriques}, Masson \&
North-Holland, 1970.

\bibitem[Gr57]{Gr57}
A.~Grothendieck, Sur quelques points d'algèbre homologique, \emph{Tôhoku
Math.\ J.}\ \textbf{9} (1957), 119--221.

\bibitem[Ja03]{Ja03}
J.~C.~Jantzen, \emph{Representations of algebraic groups}, Academic Press,
1987; 2nd~ed.\ Amer.\ Math.\ Soc., 2003.

\bibitem[GIT]{GIT}
D.~Mumford, \emph{Geometric invariant theory}, Springer-Verlag, 1965;
2nd~ed., with J.~Fogarty, 1982; 3rd~ed., with J.~Fogarty \& F.~Kirwan, 1994.

\bibitem[Ni02]{Ni02}
N.~Nitsure, Representability of~\(\GL_E\), \emph{Proc.\ Indian Acad.\ Sci.}\
\textbf{112} (2002), \No 4, 539--542.

\bibitem[Ni04]{Ni04}
N.~Nitsure, Representability of \(\Hom\) implies flatness, \emph{Proc.\
Indian Acad.\ Sci.}\ \textbf{114} (2004), \No 1, 7--14.

\bibitem[Se68]{Se68}
J.-P.~Serre, Groupes de Grothendieck des schémas en groupes réductifs
déployés, \emph{Publ.\ math.\ I.H.É.S.}\ \textbf{34} (1968), 37--52.

\bibitem[Th87]{Th87}
R.~W.~Thomason, Equivariant resolution, linearization, and Hilbert's
fourteenth problem over arbitrary base schemes, \emph{Adv.\ Maths.}\
\textbf{65} (1987), 16--34.

\end{thebibliography}
